\documentclass[utf8]{ctexart}
\usepackage{graphicx}
\usepackage{amsmath, amssymb, amsthm}
\usepackage{geometry}
\usepackage{enumitem}
\usepackage{xcolor}
\usepackage{mdframed}
\usepackage{tikz}

\usetikzlibrary{arrows.meta,positioning,calc}
\usepackage{booktabs}
\usepackage{array}

\geometry{margin=2cm}

\definecolor{critical}{RGB}{200,40,40}
\definecolor{nodeblue}{RGB}{52,120,200}
\definecolor{nodegray}{RGB}{215,228,242}
\definecolor{critnodebg}{RGB}{255,218,218}

\usetikzlibrary{graphs, positioning, arrows.meta}

\geometry{a4paper, margin=2.5cm}

% 定理环境
\newtheoremstyle{mystyle}
  {8pt}{8pt}{\kaishu}{0pt}{\bfseries}{}{1em}{}
\theoremstyle{mystyle}
\newtheorem{theorem}{定理}
\newtheorem{lemma}{引理}

% 好看的题目框
\newmdenv[
  linecolor=black!40,
  linewidth=1pt,
  backgroundcolor=gray!8,
  roundcorner=5pt,
  innertopmargin=8pt,
  innerbottommargin=8pt,
]{problembox}

\title{\textbf{离散数学2\quad 作业14}}
\date{\today}
\author{李子嘉 2024012325}
\begin{document}
\maketitle

\section*{题目1}

\begin{problembox}
 求剩余类加群 $(\mathbb{Z}_{100}, +)$ 的所有子群。
\end{problembox}
\textbf{解：} 
由于 $\mathbb{Z}_{100}$ 是一个循环群，其生成元为 $1$，因此 $\mathbb{Z}_{100}$ 的子群也都是循环群。根据循环群的性质，$\mathbb{Z}_{100}$ 的子群对应于其阶 $100$ 的所有约数。$100$ 的约数有 $1, 2, 4, 5, 10, 20, 25, 50, 100$。因此，
$\mathbb{Z}_{100}$ 的子群为 

$\{0\}, \{0, 50\}, \{0, 25, 50, 75\}, \{0, 20, 40, 60, 80\}, \{0, 10, 20, \ldots, 90\}, $

$\{0, 5, 10, \ldots, 95\}, \{0, 4, 8, \ldots, 96\}, \{0, 2, 4, \ldots, 98\}, \mathbb{Z}_{100}$ 。

\section*{题目2}
\begin{problembox}
在$S_6$中设：
$\sigma = \begin{bmatrix}
1 & 2 & 3 & 4 & 5 & 6\\
4 & 3 & 5 & 6 & 1 & 2
\end{bmatrix}$，
$\tau = \begin{bmatrix}1 & 2 & 3 & 4 & 5 & 6\\
2 & 1 & 5 & 6 & 3 & 4
\end{bmatrix}$。

试计算 $\sigma\tau, \tau\sigma, \sigma^{-1}, \sigma\tau\sigma^{-1}$，同时将它们表成对换之积。
\end{problembox}
\textbf{解：}
首先，我们将 $\sigma$ 和 $\tau$ 转换成循环表示：
$\sigma = (4\ 6\ 2\ 3\ 5\ 1)$，$\tau = (1\ 2)(3\ 5)(4\ 6)$。

随后，我们计算 $\sigma\tau$ 和 $\tau\sigma$：
$\sigma\tau = (1\ 3)(2\ 4)$，$\tau\sigma = (1\ 6)(2\ 5)$，$\sigma^{-1} = (1\ 5\ 3\ 2\ 6\ 4) = (1\ 5)(5\ 3)(3\ 2)(2\ 6)(6\ 4)$，$\sigma\tau\sigma^{-1} = (4\ 3)(5\ 1)(6\ 2)$。
\section*{题目3}
\begin{problembox}
求交错群 $A_4$ 。
\end{problembox}
\textbf{解：}
交错群 $A_4$ 是对称群 $S_4$ 的一个子群，包含所有偶置换。$A_4$ 的元素包括：
\begin{enumerate}
  \item 恒等置换 $e$。
  \item 3-循环 $(1\ 2\ 3), (1\ 3\ 2), (1\ 2\ 4), (1\ 4\ 2), (1\ 3\ 4), (1\ 4\ 3), (2\ 3\ 4), (2\ 4\ 3)$
  \item 2+2轮换 $(1\ 2)(3\ 4), (1\ 3)(2\ 4), (1\ 4)(2\ 3)$。
\end{enumerate}
因此，$A_4$ 的元素为 $e, (1\ 2\ 3), (1\ 3\ 2), (1\ 2\ 4), (1\ 4\ 2), (1\ 3\ 4), (1\ 4\ 3), (2\ 3\ 4), (2\ 4\ 3), (1\ 2)(3\ 4), (1\ 3)(2\ 4), (1\ 4)(2\ 3)$，共12个元素。
\end{document}